\documentclass[11pt]{article}
\usepackage[margin=1.05in]{geometry}
\usepackage{booktabs}
\usepackage{amsmath,amssymb}
\usepackage{microtype}
\usepackage{caption}
\usepackage{natbib}
\usepackage[colorlinks=true,linkcolor=black,citecolor=black,urlcolor=blue]{hyperref}
\captionsetup{font=small,labelfont=bf}
\bibliographystyle{plainnat}

\title{\vspace{-2.2em}The Ceiling on Skill in Head-to-Head Fantasy Football}
\author{A variance budget for a fourteen-team league, 2016--2025}
\date{September 2026}

\begin{document}
\maketitle
\vspace{-1.2em}

\begin{abstract}
\noindent We decompose weekly team scores in a fourteen-team fantasy football
league into persistent manager quality and week-to-week noise, and derive from
that decomposition the probability that the better manager wins a given matchup.
Across 770 manager-weeks the intraclass correlation is
$\widehat{\rho} = 0.039$, 95\% CI $[0.003,\,0.065]$: roughly
\emph{ninety-six percent} of the variation in a weekly score is within-manager
noise. The implied skill-to-noise ratio is $\kappa = 0.143$, so a manager one
standard deviation above average wins a given week with probability $0.557$ and
finishes a fourteen-game season above .500 with probability $0.566$. Measuring a
manager to a reliability of $0.80$ would take about ninety-nine weeks, or seven
seasons. Substituting each manager's \emph{optimal} lineup for his actual one
leaves the decomposition essentially unchanged ($\rho = 0.038$), implying that
lineup selection contributes almost nothing to the between-manager variance.
Two decisions do carry measurable weight --- avoiding early-round draft busts,
and waiver claim volume --- and we quantify both. We take the modest position
that these results describe one league over four seasons and are consistent
with, rather than proof of, a general ceiling.
\end{abstract}

\section{Introduction}

Ask a group of experienced fantasy football managers how much of the game is
skill and you will get two incompatible answers, often from the same person. One
is that the game is mostly luck. The other is that some managers make the
playoffs every year and this cannot be coincidence. Both claims are usually
offered without a number attached.

The two are in fact compatible, and the reconciliation is quantitative. If
persistent differences in manager quality are small relative to week-to-week
scoring noise, then a real and stable talent advantage will still produce a win
rate close to a coin flip in any single week, and a season of fourteen games
will be far too short to reveal it. The question is not whether skill exists but
what its \emph{magnitude} is on the scale that decides matchups.

This paper estimates that magnitude for one league. Our contributions are:

\begin{enumerate}
\item A variance decomposition of weekly team scores into persistent manager
quality and within-manager noise, with bootstrap intervals
(Section~\ref{sec:budget}).
\item A closed-form map from that decomposition to the probability that the
better manager wins, and hence to the distribution of season records
(Section~\ref{sec:winprob}).
\item An estimate of the observation horizon required to measure manager quality
at conventional reliability (Section~\ref{sec:horizon}).
\item A decomposition of the three decisions a manager controls --- the draft,
in-season acquisition, and lineup selection --- showing which carry weight and
which do not (Section~\ref{sec:decisions}).
\item An honest validation exercise in which a record-based estimate of the
talent spread is compared against the score-based one, and does not fully agree
(Section~\ref{sec:validation}).
\end{enumerate}

\subsection*{Related work}

\citet{beckersun2016} formulate fantasy football draft and lineup management as
a mixed-integer program and propose a bilinear model for weekly scoring in which
a player's expected production against a given opponent is the product of an
innate ability $u^s_i$ for statistic $s$ and the opponent's innate ability to
defend it, $w^{Ds}_j$, fitted by alternating least squares. Two of their results
bear directly on the present paper. First, their Table~3 reports mean absolute
weekly forecast errors of $6.77$ at running back against a mean of $7.46$, and
$4.61$ at tight end against a mean of $4.72$; that is, the typical miss is
between 87\% and 98\% of the typical score. Second, in a simulation in which
their optimizer is granted \emph{perfect} knowledge of the season's outcomes, it
wins a ten-team league only 24.7\% of the time against a 10\% baseline.
Omniscience is worth a factor of 2.5, not a factor of ten.

\citet{lutz2015} applies support vector regression and neural networks to six
seasons of quarterback data and reports a best mean absolute error of $6.22$
fantasy points and a mean relative error of $0.42$. His machine-learning
approaches do not materially outperform the linear baselines he compares
against.

The tier lists published at \texttt{borischen.co} are widely used in draft
preparation and are generated by model-based clustering (Gaussian mixtures via
\texttt{mclust}) over FantasyPros expert consensus rankings.\footnote{Source at
\url{https://github.com/borisachen/fftiers}.} We note that these tiers cluster
expert \emph{opinion}, not realized outcomes, and that the repository publishes
no validation of whether tier boundaries carry predictive content.

Our contribution is complementary to all three. Where \citet{beckersun2016} and
\citet{lutz2015} ask how well weekly scores can be predicted, we ask what the
irreducible residual implies for the attainable edge, and we work backwards from
outcomes rather than forwards from projections.

\section{Data}\label{sec:data}

The league is a fourteen-team, half-PPR, head-to-head league with a starting
lineup of QB / RB / RB / WR / WR / TE / FLEX (RB, WR, TE) / WRRB\_FLEX (RB, WR)
/ K / DEF and four bench slots. It has run since 2016, on MyFantasyLeague from
2016 to 2019, on Yahoo in 2020--21, and on Sleeper since 2022.

All primary analysis uses the Sleeper era, 2022--2025, obtained from the public
Sleeper API: every draft pick, every weekly matchup with the full starter list
and per-player point totals, every roster, and every transaction including
\emph{failed} waiver claims with their bid amounts. This yields 770
manager-weeks across 14 managers, 798 draft picks, and 1{,}089 reconstructable
waiver auctions.

Analysis is restricted to weeks 1--14, the fantasy regular season. Playoff weeks
are excluded because they are not played by all managers and decide nothing for
a team already eliminated.

The earlier eras are used only in Section~\ref{sec:persistence}. The 2016--17
MyFantasyLeague drafts are recoverable; 2018 returns picks with an empty player
field, 2019 returns no draft results, and the Yahoo export for 2020--21 contains
standings and matchups but no player-level or draft data. The MFL seasons are
also structurally different --- twenty franchises drafting as two
\emph{independent} ten-team drafts from a shared player pool, fifteen rounds, no
flex --- and MFL stored no scoring rules, so player values there are computed in
half-PPR and may not match what that league actually scored.

\section{Methods}

\subsection{Variance components}\label{sec:model}

Let $Y_{mt}$ be the points scored by the starting lineup of manager $m$ in week
$t$. We posit the one-way random effects model
\begin{equation}
Y_{mt} = \mu + \alpha_m + \varepsilon_{mt}, \qquad
\alpha_m \sim (0, \tau^2), \quad
\varepsilon_{mt} \sim (0, \sigma^2),
\label{eq:model}
\end{equation}
with $\alpha_m$ independent of $\varepsilon_{mt}$. Here $\alpha_m$ is the
persistent component of manager $m$'s scoring --- everything that carries from
week to week, including roster quality inherited from the draft --- and
$\varepsilon_{mt}$ is the week-specific deviation. The intraclass correlation
\begin{equation}
\rho = \frac{\tau^2}{\tau^2 + \sigma^2}
\end{equation}
is the share of variance in a single weekly score attributable to who the
manager is.

We deliberately do \emph{not} interpret $\alpha_m$ as ``skill'' in a causal
sense. It absorbs anything persistent, including a lucky draft whose value
persists all season. It is an upper bound on skill, not a measurement of it.

\subsection{From variance to win probability}\label{sec:winprob}

A head-to-head matchup is decided by the margin
\begin{equation}
D_{mm't} = Y_{mt} - Y_{m't}
= \underbrace{(\alpha_m - \alpha_{m'})}_{\textstyle \delta}
+ \underbrace{(\varepsilon_{mt} - \varepsilon_{m't})}_{\textstyle \text{noise}} .
\end{equation}
The talent gap $\delta$ enters once; the noise enters \emph{twice}, with
variance $2\sigma^2$. Assuming approximate normality of the residuals,
\begin{equation}
\Pr(m \text{ beats } m' \mid \delta)
= \Phi\!\left( \frac{\delta}{\sigma\sqrt{2}} \right).
\label{eq:pwin}
\end{equation}

It is convenient to express a manager's talent in units of the between-manager
standard deviation. Writing $\delta = c\,\tau$ for a manager $c$ standard
deviations above average facing an average opponent, and defining the
\textbf{skill-to-noise ratio}
\begin{equation}
\kappa \;\equiv\; \frac{\tau}{\sigma\sqrt{2}}
\;=\; \sqrt{\frac{\rho}{2(1-\rho)}},
\label{eq:kappa}
\end{equation}
equation~\eqref{eq:pwin} becomes simply
\begin{equation}
\Pr(\text{win a given week}) = \Phi(c\,\kappa).
\label{eq:pk}
\end{equation}
The second equality in~\eqref{eq:kappa} follows by substituting
$\tau^2 = \rho(\tau^2+\sigma^2)$ and $\sigma^2 = (1-\rho)(\tau^2+\sigma^2)$, and
shows that $\kappa$ is a function of the ICC alone. Everything downstream ---
win rates, season records, the observation horizon --- follows from $\rho$.

Season records follow immediately. With $n$ independent matchups and per-week
win probability $p = \Phi(c\kappa)$, the win total is $W \sim \text{Bin}(n, p)$
and
\begin{equation}
\Pr(\text{winning record}) = \sum_{w > n/2} \binom{n}{w} p^w (1-p)^{n-w}.
\end{equation}

\subsection{Estimation}\label{sec:estimation}

The panel is unbalanced (managers enter and leave), so we estimate
\eqref{eq:model} by unbalanced one-way ANOVA \citep{searle1992}. With $k$
managers, $n_i$ observations for manager $i$, and $N = \sum_i n_i$,
\begin{equation}
\widehat{\sigma}^2 = \mathrm{MS}_W, \qquad
\widehat{\tau}^2 = \max\!\left\{ \frac{\mathrm{MS}_B - \mathrm{MS}_W}{n_0},\; 0 \right\},
\qquad
n_0 = \frac{1}{k-1}\left( N - \frac{\sum_i n_i^2}{N} \right).
\end{equation}
The correction by $n_0$ matters: the naive between-group variance of the
manager means was $26.2$, of which $\widehat\sigma^2 / \bar n = 9.5$ is sampling
noise in the means themselves. Failing to subtract it would overstate the
persistent component by roughly a third.

Interval estimates come from a nonparametric bootstrap resampling
\emph{managers} rather than manager-weeks, since observations are clustered
within manager; $B = 2000$ replications.

\section{Results}

\subsection{The variance budget}\label{sec:budget}

Table~\ref{tab:budget} gives the decomposition.

\begin{table}[h]
\centering\small
\begin{tabular}{lrlrl}
\toprule
& \multicolumn{2}{c}{actual lineup} & \multicolumn{2}{c}{optimal lineup} \\
\cmidrule(lr){2-3}\cmidrule(lr){4-5}
& estimate & 95\% CI & estimate & 95\% CI \\
\midrule
$\sigma^2$ \ (within manager, week to week) & $499.5$ & $[450,\,550]$ & $520.0$ & $[480,\,561]$ \\
$\tau^2$ \ (between managers)               & $20.3$  & $[1.2,\,36.9]$ & $20.6$ & $[0.0,\,41.7]$ \\
$\rho$ \ (ICC)                              & $0.0390$ & $[0.0026,\,0.0653]$ & $0.0380$ & $[0.0000,\,0.0710]$ \\
$\kappa$ \ (skill-to-noise)                 & $0.1425$ & $[0.0358,\,0.1869]$ & $0.1406$ & $[0.0000,\,0.1955]$ \\
$\Phi(\kappa)$ \ (P(+1\,SD wins a week))    & $0.557$ & $[0.514,\,0.574]$ & $0.556$ & $[0.500,\,0.578]$ \\
\bottomrule
\end{tabular}
\caption{Variance budget for weekly team scores, $k=14$ managers, 770
manager-weeks, 2022--2025. ``Optimal lineup'' replaces each manager's actual
starters with the best legal lineup available from his roster that week.}
\label{tab:budget}
\end{table}

Two things stand out. First, $\rho = 0.039$: about \textbf{96\% of the variation
in a weekly score is within-manager noise.} In natural units,
$\widehat\sigma = 22.4$ points against $\widehat\tau = 4.5$ points, and the
matchup margin carries a noise standard deviation of
$\sigma\sqrt2 = 31.6$ points against a talent gap standard deviation of
$\tau\sqrt2 = 6.4$ points.

Second, and less expected: replacing every manager's actual lineup with his
\emph{optimal} lineup leaves the decomposition almost exactly where it was
($\rho$ moves from $0.039$ to $0.038$). Managers do differ in lineup efficiency
--- the observed range is $91.1\%$ to $94.1\%$ of optimal --- but the differences
are too small, relative to $\sigma$, to register in the between-manager variance.
Perfect start/sit decisions for everyone would not change who wins this league.

\subsection{Win probability and season records}

Applying~\eqref{eq:pk} and the binomial:

\begin{table}[h]
\centering\small
\begin{tabular}{lrrr}
\toprule
talent edge & $\Pr(\text{win a week})$ & $\mathbb{E}[\text{wins in 14}]$ & $\Pr(\text{record} > .500)$ \\
\midrule
$+0.5$ SD & $0.528$ & $7.40$ & $0.480$ \\
$+1$ SD   & $0.557$ & $7.79$ & $0.566$ \\
$+2$ SD   & $0.612$ & $8.57$ & $0.725$ \\
$+3$ SD   & $0.665$ & $9.32$ & $0.848$ \\
\bottomrule
\end{tabular}
\caption{Implied win probabilities. A manager three standard deviations above
average --- a figure with no counterpart in the observed data --- still loses a
third of his weeks.}
\label{tab:pwin}
\end{table}

The predicted win rate for a genuinely good manager, $55.7\%$, is worth
comparing against folk estimates. A widely-upvoted claim on
\texttt{r/fantasyfootball} from a manager reporting his own long-running
league's records is that ``bad'' owners win at 40\% and ``good'' owners at
55--60\%. Our figure is derived from the variance decomposition without
reference to any observed win rate, and lands inside that interval. It is also
consistent with \citet{beckersun2016}: an optimizer with perfect foresight won
their ten-team simulation 24.7\% of the time, against a 10\% baseline --- a
2.5-fold edge, of the same modest order as ours.

\subsection{How long it takes to measure a manager}\label{sec:horizon}

The reliability of a $k$-week mean is given by the Spearman--Brown relation
$\rho_k = k\rho / (1 + (k-1)\rho)$:

\begin{center}\small
\begin{tabular}{lrrrrr}
\toprule
weeks $k$ & 14 & 28 & 56 & 104 & 208 \\
seasons & 1.0 & 2.0 & 4.0 & 7.4 & 14.9 \\
reliability $\rho_k$ & $0.362$ & $0.532$ & $0.695$ & $0.809$ & $0.894$ \\
\bottomrule
\end{tabular}
\end{center}

Reaching a reliability of $0.80$ requires roughly \textbf{99 weeks, or seven
seasons}. A single season's record has a reliability of $0.36$ --- barely better
than a coin. This is the formal statement of why season-long fantasy standings
feel arbitrary: they are.

\subsection{Where decisions actually matter}\label{sec:decisions}

The variance budget bounds the total attainable edge but says nothing about
where it lives. We decompose the three decisions a manager controls.

\paragraph{The draft.} Each pick is scored as the player's half-PPR points in
weeks 1--14, converted to points above replacement at his own position
(replacement recomputed each season by greedily filling all fourteen teams'
starting slots), minus the mean value returned by picks within $\pm 12$ of that
pick number. The between-manager spread concentrates almost entirely in the
first four rounds: the three best drafters returned $+29.3$ points per pick in
rounds 1--2 against $-28.0$ for the three worst, a gap that shrinks
monotonically to $+4.7$ versus $-4.6$ by rounds 11--14. The mechanism is bust
avoidance rather than gem-finding --- the top three busted an early pick 8.3\% of
the time against a league rate of 16.9\%. We note that the bust threshold is a
free parameter; the \emph{ordering} of managers is stable across thresholds from
$-100$ to $0$, but the magnitude of the gap is not.

\paragraph{In-season acquisition.} Sleeper retains failed waiver claims with
their bid amounts, permitting each contested player-week to be reconstructed as
a first-price sealed-bid auction. Across fourteen managers, the correlation
between claims attempted per season and waiver production is $r = +0.888$;
between price paid per successful claim and production, $r = -0.799$; and
between claim \emph{win rate} and production, $r = -0.369$. The negative sign on
win rate is the informative one: a high win rate indicates bidding only when
certain, and certainty is expensive.

Restricting to claims made in weeks 1--6, so that acquisition date does not
confound the comparison, players won for \$0 returned $75.7$ points and players
won for \$35 or more returned $78.4$ --- a $3.6\%$ difference for an unbounded
price. Points per dollar fall from $29.9$ in the \$1--4 band to $1.7$ above \$35.
Since the FAAB budget expires worthless at season end and bench slots recycle
weekly, the marginal cost of an additional low bid is close to zero, and the
optimal policy is therefore many small claims rather than few large ones.

\paragraph{Lineup selection.} As shown in Table~\ref{tab:budget}, this
contributes nothing detectable. The entire league sits between $91.1\%$ and
$94.1\%$ of the optimal lineup; one standard deviation of lineup efficiency is
about $2.8$ points per week against a matchup noise standard deviation of $31.6$.

\subsection{Persistence across eras}\label{sec:persistence}

If manager quality is a stable trait, draft performance should correlate across
time. Comparing each manager's standardised draft score in 2016--17 against
2022--25 for the twelve managers present in both, the correlation is
$r = -0.157$, 95\% CI $[-0.67,\,+0.46]$ by Fisher's $z$. Eight of twelve changed
sign; the manager with the best four-year draft record in the current era had
the worst in the earlier one.

This estimate is attenuated by measurement error in both halves. With
within-era ICC $0.215$ for season draft scores, Spearman--Brown gives
reliabilities of $0.523$ for a four-season mean and $0.354$ for a two-season
mean, so the disattenuated correlation is
$-0.157 / \sqrt{0.523 \times 0.354} = -0.365$. Correcting for attenuation moves
the estimate \emph{further} from zero in the negative direction; it cannot
rescue a claim of persistence. To have supported $r_{\text{true}} = +0.50$ we
would have needed to observe $r_{\text{obs}} = +0.215$.

We hold this result loosely. The eras differ in league size, roster structure,
and possibly scoring, and four years separate them.

\subsection{Validation, and where it fails}\label{sec:validation}

The model in Section~\ref{sec:model} is fitted to weekly \emph{scores}. It makes
a testable prediction about season \emph{records}, which are a different
measurement. If season win rates are $\text{Bin}(14, p_m)/14$ with
$p_m = \Phi(c_m \kappa)$, then their variance decomposes into a binomial part
and a talent part. Using the delta method, the talent-induced standard deviation
of the win rate is approximately $2\phi(0)\kappa$.

\begin{center}\small
\begin{tabular}{lr}
\toprule
observed SD of season win rate & $0.1517$ \\
SD implied by a pure coin flip, $n=14$ & $0.1336$ \\
residual attributable to talent, $\sqrt{0.1517^2 - 0.1336^2}$ & $0.0718$ \\
predicted from the weekly model, $2\phi(0)\kappa$ & $0.1137$ \\
\bottomrule
\end{tabular}
\end{center}

\textbf{The two do not agree.} The weekly-score model predicts a talent spread
in season records about $1.6$ times larger than the records themselves display.
The discrepancy is not fatal --- $\widehat\tau^2$ has a 95\% interval of
$[1.2,\,36.9]$, and the record-based figure of $0.0718$ corresponds to a $\tau$
near the lower end of that interval --- but it means the point estimate in
Table~\ref{tab:budget} should be read as an \emph{upper} bound. Both
measurements agree that the talent effect is small; they disagree by a factor of
1.6 on how small, and we cannot resolve that with four seasons. Candidate
explanations include non-normality in the margin (heavier tails reduce the
leverage of a mean shift) and correlation in opponent quality induced by the
schedule.

\section{Limitations}

\paragraph{One league.} Everything here describes fourteen managers over four
seasons. The variance budget may differ in shallower leagues, where replacement
level is higher and roster differences larger.

\paragraph{$\alpha_m$ is not skill.} It is everything persistent, including a
draft that broke well. Attributing it to decision quality is an over-reading.

\paragraph{Non-independence.} The same managers recur across seasons and face
each other; matchups within a week are not independent draws. The bootstrap
clusters on managers, which addresses part but not all of this.

\paragraph{Normality.} Equation~\eqref{eq:pwin} assumes an approximately normal
margin. Weekly team scores are right-skewed sums of skewed player scores, and
the validation failure in Section~\ref{sec:validation} is consistent with this
assumption being imperfect.

\paragraph{Free parameters.} The $-50$ bust threshold and the $\pm 12$ window in
the slot-expectation curve were chosen without formal justification. The slot
curve is also fitted on the same picks it scores, inducing mild leakage; a
leave-one-out or isotonic fit would be preferable.

\section{What a manager should conclude}

The results support three practical statements and refute a fourth.

First, \textbf{the attainable edge is real but small}, on the order of five to
six percentage points of weekly win probability for a genuinely good manager.
Over fourteen games this is worth about $0.8$ wins. It is not nothing --- it is
roughly the difference between making and missing the playoffs --- but it will
not feel like control.

Second, \textbf{what edge exists is concentrated in the first four rounds of the
draft and in waiver volume}. Those are the two decisions with magnitude. The
draft evidence is the weaker of the two, resting on sixteen early picks per
manager; the waiver evidence rests on 1{,}089 auctions and is considerably more
solid.

Third, \textbf{the correct waiver policy is many small claims}, because FAAB is
a use-it-or-lose-it budget, bench slots recycle, and price buys almost nothing:
a free pickup returned $75.7$ points against $78.4$ for a \$35 claim.

Fourth, and contrary to a great deal of received advice, \textbf{weekly start/sit
optimisation is not where games are won}. It contributes nothing measurable to
between-manager variance. A manager who never watches a football game and simply
starts the higher projection while covering bye weeks gives up very little.

The honest summary is that this is a game in which a good manager should expect
to be right slightly more than half the time, and in which a single season
cannot distinguish him from a lucky one. That is not a reason to stop playing.
It is a reason to stop treating the standings as a verdict.

\begin{thebibliography}{9}
\bibitem[Becker and Sun(2016)]{beckersun2016}
Becker, A. and X.~A. Sun (2016).
\newblock An analytical approach for fantasy football draft and lineup management.
\newblock \emph{Journal of Quantitative Analysis in Sports} 12(1), 17--30.

\bibitem[Lutz(2015)]{lutz2015}
Lutz, R. (2015).
\newblock Fantasy football prediction.
\newblock \emph{arXiv preprint} arXiv:1505.06918.

\bibitem[Searle et al.(1992)]{searle1992}
Searle, S.~R., G.~Casella, and C.~E. McCulloch (1992).
\newblock \emph{Variance Components}.
\newblock Wiley, New York.
\end{thebibliography}

\end{document}
